\documentclass[11pt,a4paper]{article}
\usepackage[T1]{fontenc}
\usepackage[utf8]{inputenc}
\usepackage{lmodern,microtype}
\usepackage[margin=23mm]{geometry}
\usepackage{amsmath,amssymb,amsthm,mathtools,booktabs}
\usepackage[hidelinks]{hyperref}
\hypersetup{
  pdftitle={Exact Braid Algebra and Diamond-Norm Bounds for a Phenomenological Noisy Fibonacci-Anyon Qubit},
  pdfauthor={Andrei Preece and Boris Batenin},
  pdfsubject={Fibonacci anyons, braid relations, flagged noise, diamond-norm bounds},
  pdfkeywords={Fibonacci anyons, braid algebra, diamond norm, Lean, quantum channels}
}
\usepackage{xcolor}
\usepackage{enumitem}
\usepackage{array}
\usepackage{fancyhdr}
\pagestyle{fancy}
\fancyhf{}
\lhead{Fibonacci anyons: braid algebra and noise bounds}
\rhead{\thepage}
\setlength{\headheight}{14pt}
\setlist{nosep}
\newtheorem{theorem}{Theorem}[section]
\newtheorem{lemma}[theorem]{Lemma}
\newtheorem{proposition}[theorem]{Proposition}
\theoremstyle{remark}
\newtheorem{remark}[theorem]{Remark}
\newcommand{\Id}{\mathop{\mathrm{id}}}
\newcommand{\one}{\mathbf{1}}
\newcommand{\Cd}{\mathbb{C}}
\newcommand{\R}{\mathbb{R}}
\newcommand{\Dn}{\|\cdot\|_\diamond}
\title{\textbf{Exact Braid Algebra and Diamond-Norm Bounds for a Phenomenological Noisy Fibonacci-Anyon Qubit}}
\author{Andrei Preece and Boris Batenin\\[4pt]Independent Researchers\\Berlin, 2026}
\date{}
\begin{document}
\maketitle
\begin{abstract}
We give a finite-dimensional, reproducible model of three Fibonacci anyons constrained to total charge $\tau$. The braid operators are represented in the two-dimensional fusion-path basis using the standard $F$- and $R$-symbols. An explicit polynomial certificate proves the Yang--Baxter relation once two elementary scalar identities are supplied. We then define a physically motivated \emph{effective} quantum-noise model incorporating fusion-basis phase flips and flagged erasure from the logical space, and prove an \emph{exact} single-gate completely bounded trace-norm formula
$\|\mathcal E_{q,\ell,U}-\mathcal V_U\|_\diamond=2[\ell+(1-\ell)q]$.
More generally, an exact closed form is given for arbitrary stochastic logical Pauli errors with flagged leakage. A telescoping theorem supplies circuit-level upper bounds under compatible channel composition. The accompanying Lean sources formalize the fusion rules, polynomial braid certificate, golden-ratio $F$ identity, and abstract norm-composition bounds; the concrete complex-phase specialization and the superoperator diamond-norm instance remain proof obligations. No kernel-verification claim is made without a successful Lean build.
\end{abstract}

\section{Scope, notation, and physical assumptions}
Fibonacci anyons form a unitary braided fusion category with simple charges $\{\one,\tau\}$ and fusion rule
\begin{equation}\label{eq:fusion}
\one\otimes x=x,\qquad \tau\otimes\tau=\one\oplus\tau.
\end{equation}
The resulting fusion-path spaces grow with Fibonacci numbers, a standard setting for topological quantum computation~\cite{bonesteel2005,hormozi2007,trebst2008}. The construction below concerns a \emph{single encoded qubit}, not a microscopic anyon Hamiltonian, a complete modular tensor category in Lean, or a demonstration of hardware-level fault tolerance.

We work with linear maps between matrix algebras and take the unhalved diamond norm $\|T\|_\diamond=\|T\|_{\mathrm{cb},1}$. Thus the operationally normalized channel error is $\varepsilon(T,S)=\tfrac12\|T-S\|_\diamond$.

\section{Fusion paths and the recoupling matrix}
For three $\tau$ charges constrained to total $\tau$, the two computational basis states correspond to intermediate charges $\one$ and $\tau$. Put
\begin{equation}\label{eq:f}
\varphi=\frac{1+\sqrt5}{2},\qquad u=\varphi^{-1}=\frac{\sqrt5-1}{2},\qquad v=\varphi^{-1/2}=\sqrt{u},
\end{equation}
and define
\begin{equation}
F=\begin{pmatrix}u&v\\v&-u\end{pmatrix}.
\end{equation}
The identity $u^2+u=1$ and the choice $v^2=u$ imply $F^2=I_2$ and $F=F^\dagger$. In particular, $F$ is unitary. The labels of the two basis vectors have precise topological meaning as fusion channels; the matrix entries are not arbitrary free qubit parameters.

Let $c_n(\one),c_n(\tau)$ denote the number of paths obtained by adjoining $n$ copies of $\tau$ to the vacuum. With $(c_0(\one),c_0(\tau))=(1,0)$, fusion gives
\begin{equation}
\binom{c_{n+1}(\one)}{c_{n+1}(\tau)}=
\begin{pmatrix}0&1\\1&1\end{pmatrix}
\binom{c_n(\one)}{c_n(\tau)}.
\end{equation}
Consequently $c_3(\tau)=2$ and $c_n(\tau)$ obeys the Fibonacci recurrence.

\section{Braiding and an exact polynomial certificate}
We adopt the chiral convention
\begin{equation}\label{eq:r}
a=R_{\one}^{\tau\tau}=e^{-4\pi i/5},\qquad
b=R_{\tau}^{\tau\tau}=e^{3\pi i/5},\qquad
R=\operatorname{diag}(a,b).
\end{equation}
The opposite chirality complex-conjugates both phases. The elementary braid generators are
\begin{equation}\label{eq:braid}
B_1=R,\qquad B_2=FRF^{-1}=FRF.
\end{equation}
As products of unitary matrices, both are unitary.

\begin{theorem}[Three-strand braid relation]\label{thm:braid}
The matrices in \eqref{eq:braid} obey $B_1B_2B_1=B_2B_1B_2$.
\end{theorem}
\begin{proof}
For symbolic parameters $u,v,a,b$, set
$H=u^2+v^2-1$ and $Q=u^2(a-b)^2+ab$. Direct expansion of
$D=R(FRF)R-(FRF)R(FRF)$ yields the explicit certificate
\begin{equation}\label{eq:certificate}
D=\begin{pmatrix}
(a-b)v^2Q-a(a^2u^2+b^2v^2)H & -uv(a-b)Q-2abuv(a-b)H\\
-uv(a-b)Q-2abuv(a-b)H & -(a-b)v^2Q-b(a^2v^2+b^2u^2)H
\end{pmatrix}.
\end{equation}
We already have $H=0$. Put $z=b/a=e^{-3\pi i/5}$. Because
$2\cos(3\pi/5)=-u$, one has $z^2+uz+1=0$. From $u^2+u=1$ it follows that $u^3+2u^2=1$; hence
\[
u^2(1-z)^2+z
=u^2(z^2-2z+1)+z
=(1-u^2(u+2))z=0.
\]
Multiplication by $a^2$ gives $Q=u^2(a-b)^2+ab=0$. Substituting $H=Q=0$ in \eqref{eq:certificate} proves $D=0$.
\end{proof}

\begin{remark}
The symbolic identity \eqref{eq:certificate} is checked exactly by SymPy and encoded as a conditional theorem in \texttt{Fibonacci/BraidAlgebra.lean}. The physical-phase lemma $Q=0$, while proved above on paper, has \emph{not} yet been discharged in the supplied Lean source. The analytic claim and Lean-kernel certification therefore have different verification statuses.
\end{remark}

\section{Effective dephasing and charge-leakage model}
Let $\mathcal H_C=\Cd^2$, $\mathcal H_O=\Cd^2\oplus\Cd|L\rangle$, and $J:\mathcal H_C\hookrightarrow\mathcal H_O$ be the canonical embedding. The ideal braid channel for a unitary $U\in\{B_1,B_2\}$ (or an arbitrary braid word) is
\begin{equation}
\mathcal V_U(\rho)=JU\rho U^\dagger J^\dagger.
\end{equation}
We define two error parameters: $q\in[0,1]$ is a stochastic logical fusion-basis $Z$-flip after the braid, and $\ell\in[0,1]$ is an orthogonally flagged erasure/leakage event. With $Z=\operatorname{diag}(1,-1)$, define four $3\times2$ Kraus operators
\begin{align}\label{eq:kraus}
K_0&=\sqrt{(1-\ell)(1-q)}\,JU, &
K_1&=\sqrt{(1-\ell)q}\,JZU,\\
K_2&=\sqrt\ell\,|L\rangle\langle0|, &
K_3&=\sqrt\ell\,|L\rangle\langle1|.
\end{align}
The noisy channel is $\mathcal E_{q,\ell,U}(\rho)=\sum_{j=0}^3K_j\rho K_j^\dagger$. This is completely positive, and
\[
\sum_{j=0}^3K_j^\dagger K_j
=((1-\ell)(1-q)+(1-\ell)q+\ell)I_2=I_2,
\]
so it is trace-preserving. The orthogonal flag is a mathematical model for detected leakage from the computational fusion sector. This is \emph{phenomenological}: $q$ and $\ell$ are not inferred from a microscopic anyon bath, diffusion dynamics, or a measured hardware error rate. In particular, logical $Z$ dephasing is not claimed to be a generic local topologically protected perturbation.

\section{An exact diamond-norm expression}
\begin{theorem}[Exact one-gate error for flagged leakage and dephasing]\label{thm:exact}
For any unitary $U$ on $\mathcal H_C$ and $q,\ell\in[0,1]$, the channels above satisfy
\begin{equation}\boxed{
\|\mathcal E_{q,\ell,U}-\mathcal V_U\|_\diamond
=2\big[\ell+(1-\ell)q\big].}\label{eq:exact}
\end{equation}
Equivalently, the half-diamond error is $\varepsilon=\ell+(1-\ell)q$.
\end{theorem}
\begin{proof}
Write $t=\ell+(1-\ell)q$ and let $\mathcal V_{ZU}(\rho)=JZU\rho U^\dagger ZJ^\dagger$, $\mathcal L(\rho)=\operatorname{Tr}(\rho)|L\rangle\langle L|$. Then
\begin{equation}
\mathcal E_{q,\ell,U}-\mathcal V_U
=-t\mathcal V_U+(1-\ell)q\,\mathcal V_{ZU}+\ell\mathcal L.
\end{equation}
Each of $\mathcal V_U,\mathcal V_{ZU},\mathcal L$ is a quantum channel and therefore has diamond norm exactly one. The triangle inequality gives
$\|\mathcal E-\mathcal V_U\|_\diamond\le t+(1-\ell)q+\ell=2t$.

For a matching lower bound choose the legitimate density matrix $\rho_*=U^\dagger|+\rangle\langle+|U$, where $|\pm\rangle=(|0\rangle\pm|1\rangle)/\sqrt2$. The ideal output is $|+\rangle\langle+|$, while the noisy output is
$(1-\ell)(1-q)|+\rangle\langle+|+(1-\ell)q|-\rangle\langle-|+\ell|L\rangle\langle L|$.
In the orthonormal basis $\{|+\rangle,|-\rangle,|L\rangle\}$, their difference is diagonal with entries $-t,(1-\ell)q,\ell$. Its trace norm is $t+(1-\ell)q+\ell=2t$. Since the diamond norm dominates the induced trace norm on any particular density matrix, it is at least $2t$, completing the proof.
\end{proof}

\begin{remark}
This is an exact expression for a channel $M_2(\Cd)\to M_3(\Cd)$. It should \emph{not} be substituted uncritically into a repeated-gate circuit on $M_2(\Cd)$: after leakage, the output space has changed. A valid circuit composition must include a compatible propagation or recovery rule for the flag.
\end{remark}

\section{Extension: arbitrary stochastic Pauli noise}
The exact formula is not restricted to dephasing. Let $P_0=I$, $P_1=X$, $P_2=Y$, $P_3=Z$, and $p_\alpha\ge0$ with $\sum_{\alpha=0}^3p_\alpha=1$. Define a channel $\mathcal E_{\boldsymbol p,\ell,U}:M_2(\Cd)\to M_3(\Cd)$ by Kraus operators
\[
K_\alpha=\sqrt{(1-\ell)p_\alpha}\,JP_\alpha U\quad(0\le\alpha\le3),\qquad
K_4=\sqrt\ell\,|L\rangle\langle0|,\quad
K_5=\sqrt\ell\,|L\rangle\langle1|.
\]
\begin{theorem}[Exact Pauli-plus-erasure diamond distance]\label{thm:pauli}
For any unitary $U$, any probability vector $\boldsymbol p$, and $0\le\ell\le1$,
\begin{equation}\boxed{
\|\mathcal E_{\boldsymbol p,\ell,U}-\mathcal V_U\|_\diamond
=2\bigl[1-(1-\ell)p_0\bigr].}\label{eq:pauli}
\end{equation}
\end{theorem}
\begin{proof}
The ideal Pauli branch occurs with probability $(1-\ell)p_0$; all other branches carry nonnegative weights summing to $t=1-(1-\ell)p_0$. Subtracting $\mathcal V_U$ leaves the ideal channel with coefficient $-t$, plus CPTP Pauli branches whose weights sum to $(1-\ell)(1-p_0)$ and an erasure channel of weight $\ell$. Triangle inequality bounds the diamond norm by $2t$.

For a matching lower bound, take an auxiliary qubit and the normalized Bell state $|\Phi^+\rangle=(|00\rangle+|11\rangle)/\sqrt2$, with input $(U^\dagger\otimes I)|\Phi^+\rangle$. The four states $(JP_\alpha\otimes I)|\Phi^+\rangle$ are orthonormal; any erasure output lies in the orthogonal flagged sector. Therefore the difference between output density operators has negative eigenvalue $-t$ on the ideal Bell state, positive eigenvalues $(1-\ell)p_\alpha$ for $\alpha\ne0$, and positive flagged eigenvalues summing to $\ell$. Its trace norm is exactly $2t$.
\end{proof}
\begin{remark}
Equation~\eqref{eq:exact} is the special case $p_0=1-q$, $p_3=q$, $p_1=p_2=0$. Unlike that special case, the general Pauli lower bound uses a reference qubit. The result is a direct analytic proof, not a separately solved numerical SDP and not yet a kernel-formalized completely bounded norm theorem.
\end{remark}

\section{Compositional bounds for braid circuits}
\begin{theorem}[Telescoping diamond bound]\label{thm:telescoping}
Suppose $\mathcal E_j$ and $\mathcal U_j$ are CPTP maps with compatible input/output spaces for $j=1,\dots,n$. Then
\begin{equation}\label{eq:tel}
\|\mathcal E_n\circ\cdots\circ\mathcal E_1
-\mathcal U_n\circ\cdots\circ\mathcal U_1\|_\diamond
\le \sum_{j=1}^n\|\mathcal E_j-\mathcal U_j\|_\diamond.
\end{equation}
\end{theorem}
\begin{proof}
Insert and subtract intermediate products so that each summand differs at only a single gate. The diamond norm obeys the triangle inequality and is submultiplicative under composition. Every CPTP map has diamond norm one. Bounding each summand thus gives \eqref{eq:tel}.
\end{proof}

If compatible extended-space channels are built and each is separately proved to satisfy $\delta_j\le 2t_j$, then the operationally normalized circuit error satisfies
\begin{equation}
\varepsilon_{\mathrm{circuit}}\le\min\bigl\{1,\sum_{j=1}^n t_j\bigr\}.
\end{equation}
Any additional braid-compilation error relative to a target logical gate must be included by the triangle inequality. The results do not establish a topological threshold or a suppression law for $q$ or $\ell$ as a function of quasiparticle separation or temperature.

\section{Lean sources and scope of formal verification}
The accompanying source files implement five proof layers:
\begin{enumerate}[label=\arabic*.]
\item \texttt{Fusion.lean}: fusion outcomes, path recurrence, and the two-dimensional three-anyon fusion basis.
\item \texttt{BraidAlgebra.lean}: explicit $2\times2$ matrix multiplication and a polynomial certificate for the braid relation under $H=Q=0$.
\item \texttt{GoldenF.lean}: the golden-ratio $F$ specialization and the physical exponential phases; \emph{the exact physical-phase identity is an explicit theorem hypothesis}.
\item \texttt{QuantumNoise.lean}: explicit three-level Kraus operators for dephasing and flagged leakage; normalization of the scalar weights.
\item \texttt{DiamondComposition}, \texttt{NoiseCoefficients}, \texttt{PauliWeights}: normed-ring telescoping bounds and scalar Pauli-weight identities. A concrete mathematical instance identifying this norm with the completely bounded trace norm on superoperators is \emph{not supplied}.
\end{enumerate}

The exact quantum-channel theorems~\ref{thm:exact} and~\ref{thm:pauli} are proved above by finite-dimensional functional analysis, but neither has yet been formalized in Lean. The accompanying symbolic and numerical tests are executable using the supplied Python script. The Mac audit script requires a working installation of the pinned Lean toolchain to determine whether the Lean files type-check; no completed Lean build or kernel-axiom audit is included with these materials. The audit prints both complete theorem statements and kernel axiom dependencies; a successful build alone is not treated as proof of the intended scientific claims.

\section{Novelty and limitations}
The Fibonacci braid representation and universality are known results~\cite{bonesteel2005,hormozi2007,trebst2008}. Diamond-norm SDP methods are well-established~\cite{watrous2009,watrous2012}, and phase damping is a standard effective noise mechanism~\cite{dephasing2021}. Experimental digital demonstrations of Fibonacci-anyonic braiding exist~\cite{andersen2024}. The single-gate formula~\eqref{eq:exact} is an elementary but useful closed-form property of the explicit phenomenological flagged channel. No claim of world priority or new physical topological protection is made.

For complete machine-checked verification, further steps include: (i) kernel proof of the exact physical complex-phase identity and possibly relevant hexagon/pentagon equations; (ii) a concrete CPTP and Kraus-completeness proof for the matrix channel in Lean; (iii) a trace-norm / diamond-norm formalization yielding theorem~\ref{thm:exact} without abstract assumptions; and (iv) a fully specified composable leakage dynamics, with verified hardware calibration if physical predictive claims are intended.

\appendix
\section{Reproducibility checklist}
From the extracted project root with a network-enabled Lean 4.19.0 installation and Mathlib, run
\begin{verbatim}
python3 scripts/check_numeric.py
zsh scripts/run_lean_audit_mac.zsh
\end{verbatim}
Inspect \path{results/numerical_report.json} and \path{results/lean_axioms.log}. The second output is generated \emph{only if a real Lean run takes place}. Review \texttt{\#check} statements in \path{Fibonacci/Audit.lean}; a theorem whose conclusion is merely \texttt{True} would not certify positivity or optimality. The supplied sources contain no \texttt{sorry} placeholders, but \emph{the absence of such tokens does not certify uncompiled source}.

\begin{thebibliography}{9}
\bibitem{bonesteel2005} N. E. Bonesteel, L. Hormozi, G. Zikos, S. H. Simon, ``Braid topologies for quantum computation,'' \emph{Phys. Rev. Lett.} \textbf{95}, 140503 (2005), doi:10.1103/PhysRevLett.95.140503.
\bibitem{hormozi2007} L. Hormozi, G. Zikos, N. E. Bonesteel, S. H. Simon, ``Topological quantum compiling,'' \emph{Phys. Rev. B} \textbf{75}, 165310 (2007), doi:10.1103/PhysRevB.75.165310.
\bibitem{trebst2008} S. Trebst, M. Troyer, Z. Wang, A. W. W. Ludwig, ``A short introduction to Fibonacci anyon models,'' \emph{Prog. Theor. Phys. Suppl.} \textbf{176}, 384--407 (2008).
\bibitem{watrous2009} J. Watrous, ``Semidefinite programs for completely bounded norms,'' \emph{Theory Comput.} \textbf{5}, 217--238 (2009), doi:10.4086/toc.2009.v005a011.
\bibitem{watrous2012} J. Watrous, ``Simpler semidefinite programs for completely bounded norms,'' \emph{Chicago J. Theor. Comput. Sci.} (2013); arXiv:1207.5726.
\bibitem{dephasing2021} ``Time-varying quantum channel models for superconducting qubits,'' \emph{npj Quantum Information} (2021), doi:10.1038/s41534-021-00448-5.
\bibitem{andersen2024} ``Non-Abelian braiding of Fibonacci anyons with a superconducting processor,'' \emph{Nature Physics} \textbf{20}, 1469--1475 (2024), doi:10.1038/s41567-024-02529-6.
\end{thebibliography}
\end{document}
